\originalsection{作业4}\label{sec:ps04}
原件2页, 数学题面1页. 本组使用第\ref{ch:measure}章已明示的 Lebesgue 测度构造事实; 原题明确允许使用有限开区间并可测, 并说明相关第8讲移至第5周.
\begin{exercise}\label{ps04:1}
原题1. 对任意函数 $f:\R\to\R$, 证明 $\mathcal A=\{E\subset\R:f^{-1}(E)\text{为 Lebesgue 可测集}\}$ 是 $\sigma$ 代数.
\end{exercise}
\begin{proof}[AI 独立解答]
$f^{-1}(\R)=\R$ 可测. 原像满足 $f^{-1}(E^c)=(f^{-1}E)^c$, $f^{-1}(\bigcup E_n)=\bigcup f^{-1}(E_n)$. Lebesgue 可测集对补集和可数并封闭, 所以 $\mathcal A$ 具备全部 $\sigma$ 代数公理. 不要求 $f$ 本身可测.
\end{proof}
\begin{exercise}\label{ps04:2}
原题2. 设 $E\subset\R$, $m^*(E)<\infty$. 证明 $E$ 可测当且仅当每个 $\eps>0$ 都存在有限个开区间的并 $U$, 使 $m^*(U\triangle E)<\eps$. 这是 Littlewood 第一原则. 原提示: 逆向证明 $m^*(A\cap E)+m^*(A\cap E^c)\le m^*(A)+\eps$ 对任意 $A\subset\R$ 成立.
\end{exercise}
\begin{proof}[AI 独立解答]
正向: 由外测度定义取开区间覆盖 $E\subset O=\bigcup_{j\ge1}I_j$, 使 $\sum|I_j|<m(E)+\eps/2$. $O$ 可测且有限测度, $m(O\setminus E)=m(O)-m(E)<\eps/2$. 取 $N$ 使 $\sum_{j>N}|I_j|<\eps/2$, 令 $U=\bigcup_{j\le N}I_j$. 则 $U\setminus E\subset O\setminus E$, $E\setminus U\subset\bigcup_{j>N}I_j$, 故对称差的外测度小于 $\eps$.

逆向: 固定 $A$. 若 $m^*(A)=\infty$, Carathéodory 的所需方向自动成立. 否则取 $m^*(U\triangle E)<\delta$. 因 $U$ 可测,
\begin{align*}
m^*(A\cap E)&\le m^*(A\cap U)+\delta,\\
m^*(A\setminus E)&\le m^*(A\setminus U)+\delta,\\
m^*(A\cap E)+m^*(A\setminus E)&\le m^*(A)+2\delta.
\end{align*}
令 $\delta\downarrow0$ 得分割不等式; 另一方向由外测度的次可加性得到. 故 $E$ 满足可测性定义. 取 $\delta=\eps/2$ 即为原提示的不等式.
\end{proof}
\begin{exercise}\label{ps04:3}
原题3. 设 $E$ 可测. (a) 对每个 $x\in\R$, 证明 $E+x$ 可测. (b) 对每个 $r>0$, 证明 $rE=\{ry:y\in E\}$ 可测.
\end{exercise}
\begin{proof}[AI 独立解答]
(a) 对任意 $A\subset\R$, 将 $A-x$ 由 $E$ 分割, 再用题\ref{ps03:3}的平移不变性, 得
$m^*(A)=m^*(A\cap(E+x))+m^*(A\setminus(E+x))$, 故可测.

(b) 对区间覆盖缩放, 再反向用 $r^{-1}$, 得 $m^*(rB)=r\,m^*(B)$ 对任意 $B$ 成立. 将 $r^{-1}A$ 由 $E$ 分割后乘 $r$, 得同一可测性等式. $r>0$ 使扩展非负实数的乘法不出现未定义运算.
\end{proof}
